\documentclass[11pt,a4paper]{article}
\usepackage[utf8]{inputenc}
\usepackage[T1]{fontenc}
\usepackage{lmodern}
\usepackage[brazil]{babel}
\usepackage{amsmath,amssymb,amsthm}
\usepackage{geometry}
\usepackage{hyperref}
\usepackage{booktabs}
\usepackage{enumitem}
\usepackage{xcolor}

\geometry{margin=2.5cm}

\definecolor{remark}{rgb}{0.12,0.37,0.55}

\newtheorem{theorem}{Proposição}
\newcommand{\E}{\mathbb{E}}
\newcommand{\dd}{\mathrm{d}}
\newcommand{\hx}{\hat{x}}
\newcommand{\hpi}{\hat{\pi}}

\title{\textbf{Capítulo 7 --- Modelos DSGE Novo-Keynesianos}\\[4pt]
       {\large Derivações do Modelo Canônico NK}\\[2pt]
       {\normalsize PPGEco-CCJE-UFES · Macroeconomia I · Romer (2019) Cap.\ 7}}
\author{Projeto de Estudo --- \textit{macro-romer}}
\date{2026}

\begin{document}
\maketitle
\tableofcontents
\bigskip

% ============================================================
\section{Modelo NK canônico --- três equações}
% ============================================================

O modelo Novo-Keynesiano canônico consiste em três equações log-linearizadas
em torno do estado estacionário:

\paragraph{IS dinâmica (equação de Euler das famílias):}
\begin{equation}\label{eq:is}
  x_t = \E_t[x_{t+1}] - \frac{1}{\sigma}(i_t - \E_t[\pi_{t+1}] - r^n),
\end{equation}
onde $x_t$ é o hiato do produto, $i_t$ é a taxa nominal de juros,
$\pi_t$ é a inflação e $r^n = -\ln\beta$ é a taxa real natural.

\paragraph{Curva de Phillips Novo-Keynesiana (CPNK):}
\begin{equation}\label{eq:nkpc}
  \pi_t = \beta\,\E_t[\pi_{t+1}] + \kappa\, x_t + u_t,
\end{equation}
onde $u_t$ é um choque de custo-push (AR1 com persistência $\rho_u$).

\paragraph{Regra de Taylor:}
\begin{equation}\label{eq:taylor}
  i_t = r^n + \phi_\pi \pi_t + \phi_x x_t + v_t,
\end{equation}
onde $v_t$ é um choque de política monetária (AR1 com persistência $\rho_v$).

% ============================================================
\section{Estado estacionário e taxa natural}
% ============================================================

No estado estacionário sem choques ($v_t = u_t = 0$):
\[
  x^* = 0,\qquad \pi^* = 0,\qquad i^* = r^n.
\]

A taxa real natural é:
\begin{equation}
  r^n = -\ln\beta \approx \frac{1-\beta}{\beta}.
\end{equation}

Para $\beta = 0{,}99$ (trimestral), $r^n \approx 0{,}01$ por período,
ou cerca de $4\%$ ao ano.

% ============================================================
\section{Solução por coeficientes indeterminados}
% ============================================================

Substituindo \eqref{eq:taylor} em \eqref{eq:is}:
\begin{equation}\label{eq:is_taylor}
  \left(1 + \frac{\phi_x}{\sigma}\right) x_t
  + \frac{\phi_\pi}{\sigma}\,\pi_t
  = \E_t[x_{t+1}] + \frac{1}{\sigma}\,\E_t[\pi_{t+1}] - \frac{v_t}{\sigma}.
\end{equation}

\subsection{Choque de política monetária ($v_t$)}

Suponha que $u_t = 0$ e $v_t = \rho_v\,v_{t-1} + \varepsilon_t^v$.
Adivinha-se a solução linear:
\[
  x_t = \psi_{xv}\,v_t, \qquad \pi_t = \psi_{\pi v}\,v_t.
\]

\paragraph{Da CPNK \eqref{eq:nkpc}:}
\[
  \psi_{\pi v}\,v_t = \beta\,\rho_v\,\psi_{\pi v}\,v_t + \kappa\,\psi_{xv}\,v_t
  \;\Longrightarrow\;
  \psi_{\pi v}(1 - \beta\rho_v) = \kappa\,\psi_{xv}.
\]
Portanto:
\begin{equation}\label{eq:psi_piv}
  \psi_{\pi v} = \frac{\kappa\,\psi_{xv}}{1-\beta\rho_v}.
\end{equation}

\paragraph{Da equação IS modificada \eqref{eq:is_taylor}:}
Dividindo por $v_t$ e usando $\E_t[v_{t+1}] = \rho_v v_t$:
\[
  \psi_{xv}\!\left(1 + \frac{\phi_x}{\sigma} - \rho_v\right)
  + \psi_{\pi v}\!\frac{\phi_\pi - \rho_v}{\sigma} = -\frac{1}{\sigma}.
\]
Substituindo \eqref{eq:psi_piv}:
\[
  \psi_{xv}\!\left[\underbrace{(1-\rho_v)
    + \frac{1}{\sigma}\!\left(\phi_x
    + \frac{\kappa(\phi_\pi-\rho_v)}{1-\beta\rho_v}\right)}_{\Delta_v}\right]
  = -\frac{1}{\sigma}.
\]

\begin{equation}\label{eq:solution_v}
  \boxed{
    \psi_{xv} = \frac{-1}{\sigma\,\Delta_v},
    \qquad
    \psi_{\pi v} = \frac{-\kappa}{\sigma\,\Delta_v\,(1-\beta\rho_v)},
  }
\end{equation}
onde $\Delta_v = (1-\rho_v) + \bigl[\phi_x + \kappa(\phi_\pi-\rho_v)/(1-\beta\rho_v)\bigr]/\sigma > 0$
quando o princípio de Taylor é satisfeito.

\paragraph{Interpretação:}
\begin{itemize}
  \item $\psi_{xv} < 0$: aperto monetário ($v_t > 0$) reduz o hiato do produto.
  \item $\psi_{\pi v} < 0$: aperto monetário reduz a inflação.
  \item Quanto maior $\phi_\pi$ (mais agressivo), maior $|\psi_{\pi v}|/|\psi_{xv}|$
    (inflação responde mais do que o produto).
\end{itemize}

\subsection{Choque de custo-push ($u_t$)}

Agora $v_t = 0$ e $u_t = \rho_u\,u_{t-1} + \varepsilon_t^u$.
Adivinhe:
\[
  x_t = \psi_{xu}\,u_t, \qquad \pi_t = \psi_{\pi u}\,u_t.
\]

\paragraph{Da CPNK:}
\[
  \psi_{\pi u}(1-\beta\rho_u) = \kappa\,\psi_{xu} + 1
  \;\Longrightarrow\;
  \psi_{\pi u} = \frac{1 + \kappa\,\psi_{xu}}{1-\beta\rho_u}. \tag{3}
\]

\paragraph{Da IS modificada (com $v_t=0$):}
\[
  \psi_{xu}\!\left(1 + \frac{\phi_x}{\sigma} - \rho_u\right)
  + \psi_{\pi u}\!\frac{\phi_\pi - \rho_u}{\sigma} = 0.
\]
Substituindo (3):
\[
  \psi_{xu}\!\underbrace{\left[(1-\rho_u)
    + \frac{\phi_x + \kappa(\phi_\pi-\rho_u)/(1-\beta\rho_u)}{\sigma}\right]}_{D_u}
  = -\frac{\phi_\pi-\rho_u}{\sigma(1-\beta\rho_u)}.
\]

\begin{equation}\label{eq:solution_u}
  \boxed{
    \psi_{xu} = \frac{-(\phi_\pi-\rho_u)}{\sigma(1-\beta\rho_u)\,D_u},
    \qquad
    \psi_{\pi u} = \frac{1 + \kappa\,\psi_{xu}}{1-\beta\rho_u}.
  }
\end{equation}

\paragraph{Interpretação:}
\begin{itemize}
  \item $\psi_{\pi u} > 0$: choque de custo-push eleva a inflação.
  \item $\psi_{xu} < 0$: choque de custo-push contrai o produto
    (estagflação --- o dilema de política monetária).
  \item O BC não pode estabilizar $x$ e $\pi$ simultaneamente
    após um choque de custo-push.
\end{itemize}

% ============================================================
\section{Condição de Blanchard-Kahn (princípio de Taylor)}
% ============================================================

\subsection{Matriz do sistema expectacional}

Eliminando expectativas do sistema (sem choques), o equilíbrio racional
satisfaz:
\[
  \begin{pmatrix}\E_t[x_{t+1}]\\\E_t[\pi_{t+1}]\end{pmatrix}
  = M\begin{pmatrix}x_t\\\pi_t\end{pmatrix},
\]
onde:
\begin{equation}\label{eq:system_matrix}
  M = \begin{pmatrix}
    1 + \dfrac{\phi_x}{\sigma} + \dfrac{\kappa}{\beta\sigma}
    & \dfrac{\phi_\pi}{\sigma} - \dfrac{1}{\beta\sigma} \\[8pt]
    -\dfrac{\kappa}{\beta}
    & \dfrac{1}{\beta}
  \end{pmatrix}.
\end{equation}

\subsection{Condição de determinação}

Como $x_t$ e $\pi_t$ são ambas variáveis de salto (jump variables),
a condição de Blanchard-Kahn exige que os \emph{dois} valores próprios
de $M$ estejam fora do círculo unitário ($|\lambda_{1,2}| > 1$).

O determinante de $M$ é:
\[
  \det(M) = \frac{1 + (\phi_x + \kappa\phi_\pi)/\sigma}{\beta} > 1
  \quad\text{se}\quad \phi_x + \kappa\phi_\pi > \sigma(\beta-1),
\]
o que é satisfeito para $\phi_x,\phi_\pi \geq 0$ e $\beta < 1$.

A condição suficiente (aproximação analítica) é o \textbf{princípio de Taylor}:
\begin{equation}\label{eq:taylor_principle}
  \boxed{\phi_\pi > 1 \;\text{ (e }\phi_x \geq 0).}
\end{equation}

Mais precisamente, a região de determinação satisfaz:
\[
  \kappa(\phi_\pi - 1) + (1-\beta)\phi_x > 0.
\]
Para $\beta \approx 1$ (frequência trimestral), o termo $(1-\beta)\phi_x \approx 0$
e a condição reduz-se a $\phi_\pi > 1$.

% ============================================================
\section{Política monetária ótima e fronteira eficiente}
% ============================================================

\subsection{Perda da função-objetivo do banco central}

O banco central minimiza a perda quadrática esperada:
\begin{equation}
  L = \E_0\!\left[\sum_{t=0}^\infty \beta^t\!
    \left(\pi_t^2 + \lambda_x\,x_t^2\right)\right],\quad \lambda_x > 0.
\end{equation}

\subsection{Fronteira de variância (Taylor 1979)}

Variando $\phi_\pi$ (e fixando $\phi_x$), o BC traça a fronteira eficiente
entre $\mathrm{Var}(\pi_t)$ e $\mathrm{Var}(x_t)$.

\begin{itemize}
  \item $\phi_\pi \to \infty$: $\mathrm{Var}(\pi) \to 0$, mas
    $\mathrm{Var}(x) \uparrow$ (inflação totalmente estabilizada,
    mas produto volátil).
  \item $\phi_\pi \to 1^+$: $\mathrm{Var}(x) \to \min$, mas
    $\mathrm{Var}(\pi) \uparrow$ (produto estabilizado, inflação volátil).
\end{itemize}

A fronteira é implementada numericamente em \texttt{ch07\_nk.py} via o método
\texttt{policy\_frontier()}, que simula os momentos de segunda ordem para
diferentes valores de $\phi_\pi$.

% ============================================================
\section{Parâmetros de calibração}
% ============================================================

\begin{table}[h!]
  \centering
  \caption{Parâmetros do modelo NK canônico}
  \begin{tabular}{llcc}
    \toprule
    Parâmetro & Descrição & Valor base & Referência \\
    \midrule
    $\beta$    & Fator de desconto             & $0{,}99$  & Trimestral \\
    $\sigma$   & Inv.\ elasticidade intertemp. & $1{,}0$   & Log-utilidade \\
    $\kappa$   & Slope da CPNK                 & $0{,}1$   & Galí (2008) \\
    $\phi_\pi$ & Peso da inflação (Taylor)     & $1{,}5$   & Taylor (1993) \\
    $\phi_x$   & Peso do hiato (Taylor)        & $0{,}5$   & Taylor (1993) \\
    $\rho_v$   & Persistência choque de demanda & $0{,}5$  & Estimado \\
    $\rho_u$   & Persistência choque de oferta & $0{,}5$   & Estimado \\
    \bottomrule
  \end{tabular}
\end{table}

\noindent A taxa real natural é $r^n = -\ln(0{,}99) \approx 1{,}005\%$ por trimestre
($\approx 4\%$ ao ano). O slope da CPNK ($\kappa = 0{,}1$) implica,
pela fórmula de Calvo, $\alpha \approx 0{,}75$ (firmas reajustam preços
a cada 4 trimestres em média).

% ============================================================
\section{Funções de resposta ao impulso}
% ============================================================

Dada a solução \eqref{eq:solution_v}--\eqref{eq:solution_u}, as IRFs
para os choques são:

\paragraph{Choque de demanda $v_0 > 0$, $u_t = 0$:}
\[
  v_t = \rho_v^t\,v_0,\quad
  x_t = \psi_{xv}\,v_t < 0,\quad
  \pi_t = \psi_{\pi v}\,v_t < 0,\quad
  i_t = r^n + \phi_\pi\pi_t + \phi_x x_t + v_t > 0.
\]

\paragraph{Choque de oferta $u_0 > 0$, $v_t = 0$:}
\[
  u_t = \rho_u^t\,u_0,\quad
  x_t = \psi_{xu}\,u_t < 0,\quad
  \pi_t = \psi_{\pi u}\,u_t > 0.
\]
O banco central aumenta $i_t$ para conter a inflação, mas isso aprofunda
a contração do produto --- o dilema da política monetária após choques de oferta.

\end{document}
